\section{Conclusions and outlook}
\label{sec:conclusions}

Lebed asked how often diquarks form; the dynamical answer is
\emph{more often than the static estimate says}.  In our pilot
ensembles the fraction of baryons assembled through a persistent
diquark intermediate is $0.86$--$1.00$ at every point scanned ---
expansion rates spanning a factor of four and couplings up to
$\Gamma\approx1.2$ --- against a static band of $26$--$31\%$ at the same
density: binding, once achieved, is remembered.  Species yields fall
geometrically with constituent number ($\sim60\times$ per quark from
meson to baryon at $\Gamma\approx0.8$, the $3\to4$ step relaxing from
$15\times$ to $8\times$ as the coupling grows), the coalescence-like
structure whose slope is this model's quantitative output.  The
annihilation sector behaves as designed: photon-channel depletion is
calibrated ($\tau_{\rm eff}\approx7.2\,\tau_{\rm bound}$ for compact
orbits), thermal re-injection produces net strangeness, and interaction
correlations shift the static probability by at most a few percentage
points.  On the quantum side, the four-tier hierarchy makes the same
ratios renormalization-group-safe observables of real-time lattice QCD:
integer-quantized charges and spectral projectors carry the rigor,
asymptotic pair spectra carry the broad resonances, and flowed
projectors carry --- deliberately scheme-labeled --- the
compact-versus-molecular structure question.

Outlook: Wong-equation continuous color as the principled classical
upgrade (cf.\ the strongly-coupled classical QGP program of
Ref.~\cite{Gelman:2006xw}); the 3D small-flow-time expansion for the equal-time
energy-momentum tensor as a standalone lattice-theory project; and, as
hardware matures, the fireball-quench protocol of
Sec.~\ref{sec:protocols} executed first in $1{+}1$D QCD, where every
ingredient of the hierarchy already has a small-scale analog.
